The neuroscience of curiosity has converged on a surprisingly specific picture over the last decade. Curiosity is not a single thing---it operates through at least two distinct neural mechanisms, and understanding both is essential before we can ask whether AI has anything comparable.

\subsection{The Information Gap Theory}\label{subsec:info-gap}

The dominant psychological framework, introduced by \citet{loewenstein1994}, treats curiosity as an aversive state---a felt \emph{deprivation} arising from the gap between what you know and what you want to know. This is not metaphorical. Gruber, Ranganath, and colleagues have shown through fMRI that:

\begin{itemize}[leftmargin=*]
  \item \textbf{Curiosity activates the dopaminergic reward circuit} (VTA, nucleus accumbens) \emph{during the anticipation phase}---before the answer arrives. The same circuitry that fires when you anticipate food, money, or sex fires when you are curious about a trivia answer \citep{murphy2021,rueterbories2024}.

  \item \textbf{Curiosity satisfaction triggers reward-processing ERPs} (Reward Positivity, P3, Late Positive Potential). The brain literally processes the relief of curiosity as a reward event, with the same electrophysiological signatures as receiving a monetary gain \citep{rueterbories2024}.

  \item \textbf{Curiosity enhances memory formation} through hippocampal-dopaminergic coupling. High-curiosity states do not just feel different---they produce measurably better encoding, and this effect persists over time without requiring sleep consolidation \citep{stare2018}.

  \item \textbf{The enhancement is specific, not generalized.} Contrary to early assumptions, curiosity does not create a diffuse ``sponge state'' that absorbs everything. Pre-testing effects (the memory boost from trying to answer before seeing the answer) are target-specific and do not spill over to incidental information \citep{hollins2022}.
\end{itemize}

\subsection{The Five Dimensions (5DCR Framework)}\label{subsec:5dcr}

\citet{kashdan2018} identify curiosity as having distinct subtypes in their Five-Dimensional Curiosity Scale Revised:

\begin{enumerate}[leftmargin=*]
  \item \textbf{Joyous Exploration}---delight in seeking new knowledge.
  \item \textbf{Deprivation Sensitivity}---discomfort from incomplete information (the Loewenstein gap).
  \item \textbf{Stress Tolerance}---willingness to accept uncertainty and anxiety during exploration.
  \item \textbf{Thrill Seeking}---openness to novel, intense experiences.
  \item \textbf{Social Curiosity}---interest in other people's thoughts and behaviors.
\end{enumerate}

These are not just survey categories---they correlate with different behavioral profiles, different personality structures, and (importantly for AI) different functional roles in adaptive behavior.

\subsection{The Developmental Angle}\label{subsec:developmental}

Curiosity in humans is \emph{developmental}. Children's curiosity is multimodal (cognitive, emotional, sensorimotor), context-dependent, and emerges from the affordances of their environment---objects with higher ``Exploration Quality'' (EQO) spark more curiosity, and children do not just want to know, they want to \emph{touch, smell, and manipulate} \citep{baruch2025}.

This is significant: human curiosity is embodied. It arises from the intersection of a sensing body, an environment with explorable structure, and a cognitive system that detects gaps.
